%!TEX root = ../thesis.tex

%hakon -- macro-level model chapter

\chapter{Macro-Level Model: SIRVHDB and MPC}
\label{ch:macro_model}

\section{Model Development and Rationale}
\label{sec:model_dev}

%hakon
The macro model is based on a compartmental epidemic framework in which the population is divided into discrete disease compartments.
The model was implemented in MATLAB/Simulink for modular construction and tight integration with the MPC system.
The epidemic dynamics of each city are encapsulated in a dedicated Simulink subsystem, allowing straightforward extension to additional cities by duplicating the subsystem block.
The model assumes no births and no non-disease-related deaths, keeping the total population constant when the Dead compartment is counted.

The model was built incrementally.
The baseline SIR framework was first extended by adding a Vaccinated ($V$) compartment, enabling transitions from susceptible to vaccinated at a constant daily rate $\nu$.
Breakthrough infection was then introduced following the SIRVB structure of \citet{ariyaratne2025}, allowing vaccinated individuals to be infected with reduced probability.
A Dead ($D$) compartment was added to track mortality, and a Hospitalised ($H$) compartment was introduced so that hospital load could serve as the controlled output for the MPC.
Finally, the single-city model was extended to two cities through a travel matrix for long-term population redistribution and a commuting-based mixing model for short-term cross-city contacts.

\section{SIRVHDB Compartmental Model}
\label{sec:sirvhdb}

%hakon
For each city $c \in \{1, 2\}$ the population is partitioned into seven compartments:
Susceptible ($S_c$), Infected ($I_c$), Vaccinated-Infected ($I_{v,c}$), Recovered ($R_c$), Vaccinated ($V_c$), Hospitalised ($H_c$), and Dead ($D_c$).
The total city population $N_c = S_c + I_c + I_{v,c} + R_c + V_c + H_c$ is conserved (excluding deaths).

The full single-city ODE system is:
\begin{align}
  \frac{dS_c}{dt}   &= -\lambda_c S_c - \nu_c S_c + \omega_R R_c + \omega_V V_c, \label{eq:dS} \\
  \frac{dI_c}{dt}   &= \lambda_c S_c - (\gamma + \mu + h)\,I_c, \label{eq:dI} \\
  \frac{dI_{v,c}}{dt} &= \lambda_{v,c} V_c - (\gamma_v + \mu_v + h_v)\,I_{v,c}, \label{eq:dIv} \\
  \frac{dV_c}{dt}   &= \nu_c S_c - \lambda_{v,c} V_c - \omega_V V_c, \label{eq:dV} \\
  \frac{dH_c}{dt}   &= h\,I_c + h_v\,I_{v,c} - (\gamma_h + \mu_h)\,H_c, \label{eq:dH} \\
  \frac{dR_c}{dt}   &= \gamma\,I_c + \gamma_v\,I_{v,c} + \gamma_h\,H_c - \omega_R R_c, \label{eq:dR} \\
  \frac{dD_c}{dt}   &= \mu\,I_c + \mu_v\,I_{v,c} + \mu_h\,H_c. \label{eq:dD}
\end{align}

Table~\ref{tab:parameters} summarises the model parameters and their values.

\begin{table}[ht]
\centering
\caption{SIRVHDB model parameters and baseline values.}
\label{tab:parameters}
\begin{tabular}{@{}llll@{}}
\toprule
\textbf{Parameter} & \textbf{Description} & \textbf{Value} & \textbf{Units} \\
\midrule
$\beta_{c,\mathrm{nom}}$ & Nominal transmission rate             & 0.45     & day$^{-1}$ \\
$k$             & Maximum intervention effectiveness       & 0.70     & dimensionless \\
$\gamma$        & Recovery rate, unvaccinated infected     & $1/12$   & day$^{-1}$ \\
$\gamma_v$      & Recovery rate, vaccinated breakthrough   & $1/9$    & day$^{-1}$ \\
$\gamma_h$      & Hospital discharge rate                  & $1/25$   & day$^{-1}$ \\
$h$             & Hospitalisation rate, unvaccinated       & 0.05     & day$^{-1}$ \\
$h_v$           & Hospitalisation rate, vaccinated         & 0.02     & day$^{-1}$ \\
$\mu$           & Fatality rate, unvaccinated infected     & 0.001    & day$^{-1}$ \\
$\mu_v$         & Fatality rate, vaccinated breakthrough   & 0.0005   & day$^{-1}$ \\
$\mu_h$         & Fatality rate, hospitalised              & 0.02     & day$^{-1}$ \\
$\eta$          & Relative infectiousness of $I_v$         & 0.60     & dimensionless \\
$\omega_R$      & Waning immunity, post-recovery           & $1/110$  & day$^{-1}$ \\
$\omega_V$      & Waning immunity, vaccine-induced         & $1/140$  & day$^{-1}$ \\
$\nu_1$         & Vaccination rate, City 1                 & 0.001    & day$^{-1}$ \\
$\nu_2$         & Vaccination rate, City 2                 & 0.002    & day$^{-1}$ \\
$H_{{\mathrm{cap}},c}$ & Hospital bed capacity, city $c$     & see text & agents \\
\bottomrule
\end{tabular}
\end{table}

\subsection{Force of Infection}
\label{sec:foi}

%hakon
The force of infection acting on susceptible agents in city $c$ is
\begin{equation}
  \lambda_c(t) = \beta_c(t)\,\frac{I_{c,\mathrm{eff}}(t)}{N_c},
  \label{eq:foi}
\end{equation}
where the effective infectious population mixes across city boundaries through the commuting matrix (see Section~\ref{sec:intercity}).
The force of infection on vaccinated agents is
\begin{equation}
  \lambda_{v,c}(t) = (1 - \varepsilon)\,\lambda_c(t),
  \label{eq:foiv}
\end{equation}
where $\varepsilon$ is the vaccine effectiveness against infection.

\subsection{Seasonal Transmission Forcing}
\label{sec:seasonal}

%hakon
Seasonal variation in transmission is incorporated by modulating the nominal transmission rate with a sinusoidal function of time:
\begin{equation}
  \beta_{c,\mathrm{nom}}(t) = \bar{\beta}\left(1 + 0.4\sin\!\left(\frac{2\pi t}{365}\right)\right),
\end{equation}
where $\bar{\beta} = 0.45$ day$^{-1}$ is the annual mean.
This produces higher transmission in winter months, consistent with the observed seasonality of respiratory viruses.

\section{Multi-City Interaction}
\label{sec:intercity}

%hakon
The dynamics of city $c$ are coupled to city $c'$ through two distinct mechanisms: long-term travel and commuting-based mixing.

\subsection{Travel Matrix}
\label{sec:travel}

%hakon
Long-term movement between cities is modelled using a travel matrix $T$.
The travel matrix defines the rate at which individuals move between the cities and is applied to all compartments except the Dead and Hospitalised:
\begin{equation}
  \frac{dX_c}{dt}\bigg|_{\mathrm{travel}} = \sum_{c'} T_{cc'}\,X_{c'},
\end{equation}
where $X_c$ represents any mobile compartment in city $c$.
The diagonal entries $T_{cc}$ represent net outflow from city $c$, and the off-diagonal entries $T_{cc'}$ ($c' \neq c$) represent inflow from city $c'$.
In the implementation, the travel matrix is:
\begin{equation}
  T = \begin{pmatrix} -0.2 & 0.2 \\ 0.1 & -0.1 \end{pmatrix},
\end{equation}
meaning 20\% of the population per day moves from City~1 to City~2, and 10\% from City~2 to City~1.
This results in redistribution of individuals across the cities, coupling their epidemic dynamics.

\subsection{Commuting-Based Mixing}
\label{sec:commuting}

%hakon
In addition to long-term travel, short-term interactions between the populations are modelled using a commuting-based approach.
In this mechanism, individuals do not relocate but may be exposed to infections from the other city's infectious population and return home.
The effective infectious population in city $c$ is
\begin{equation}
  I_{c,\mathrm{eff}}(t) = I_c(t) + \eta\,I_{v,c}(t),
\end{equation}
and the mixed effective infectious population seen by city $c$ is
\begin{equation}
  I_{c,\mathrm{mix}}(t) = \sum_{c'} M_{cc'}\,I_{c',\mathrm{eff}}(t),
\end{equation}
where $M$ is the commuting matrix:
\begin{equation}
  M = \begin{pmatrix} 0.85 & 0.15 \\ 0.10 & 0.90 \end{pmatrix}.
\end{equation}
The diagonal entries $M_{cc}$ represent the proportion of within-city contacts, and the off-diagonal entries $M_{cc'}$ represent cross-city contacts due to commuting.
The force of infection in city $c$ then becomes
\begin{equation}
  \lambda_c(t) = \beta_c(t)\,\frac{I_{c,\mathrm{mix}}(t)}{N_c}.
\end{equation}
This configuration allows individuals to be exposed to infections in the other city without permanently relocating there.

\section{MPC Formulation}
\label{sec:mpc_form}

%hakon
The MPC controller is designed to regulate the combined hospital load ratio across both cities.
The controlled output is the normalised hospital load:
\begin{equation}
  H_{{\mathrm{load}}}(t) = \frac{H_1(t) + H_2(t)}{H_{{\mathrm{cap}},1} + H_{{\mathrm{cap}},2}},
\end{equation}
and the reference level is $H_{\mathrm{ref}} = 0.8$, corresponding to 80\% of national hospital capacity.

\subsection{Control Input and Transmission Modulation}
\label{sec:control_input}

%hakon
The control inputs are $u_1(t)$ and $u_2(t)$, each constrained to $[0, 1]$, where $u_c = 0$ means no intervention and $u_c = 1$ means maximum restriction.
The control input modulates the transmission rate through
\begin{equation}
  \beta_c(t) = \beta_{c,\mathrm{nom}}(t)\,(1 - k\,u_c(t)),
  \label{eq:beta_control}
\end{equation}
where $k = 0.7$ is the maximum intervention effectiveness.
Thus the strongest intervention reduces transmission by 70\%.

To ensure realistic policy implementation, the control signal is passed through a rate limiter (restricting abrupt changes in intervention) and a transport delay of 5~days representing the delayed effect of policy changes on actual disease transmission.

\subsection{MPC Objective and Constraints}
\label{sec:mpc_objective}

%hakon
The MPC minimises the weighted quadratic cost function over a finite prediction horizon $N_p$:
\begin{equation}
  J = \sum_{i=1}^{N_p} w_{{\mathrm{mo}},i}\,(H_{{\mathrm{ref}}} - H_{{\mathrm{load}},i})^2 + \sum_{i=1}^{M} \omega_{u,i}\,w_{u,i}\,(\Delta u_i)^2,
  \label{eq:mpc_cost}
\end{equation}
where $w_{{\mathrm{mo}},i}$ is the output tracking weight, $w_{u,i}$ is the weight penalising rapid changes in the manipulated variable, and $\omega_{u,i}$ is the absolute control weight.
The MPC Toolbox (MATLAB R2025a) is used with the following parameters:
\begin{itemize}
  \item Prediction horizon: $N_p = 40$~days.
  \item Control horizon: $M = 10$~days.
  \item Sampling time: 1~day.
  \item Output tracking weight: 1.0.
  \item Control rate weight: 0.2.
  \item Absolute control weight: 0.05.
\end{itemize}

\subsection{Linearised Plant Model and Kalman Filter}
\label{sec:kalman}

%hakon
The MATLAB Simulink Model Linearizer was used to obtain a linearised representation of the SIRVHDB dynamics around an operating point.
This linearised model (\texttt{plant.mat}) serves as the internal prediction model for the MPC controller.
The linear approximation is necessary because the MPC Toolbox solves a quadratic program (QP) at each step, which requires a linear plant model.

The MPC Toolbox employs a Kalman filter as the default state estimator.
The Kalman filter estimates unmeasured compartment states from the observed hospital load output $H_{\mathrm{load}}(t)$, enabling the controller to update its predictions based on new measurements at each sample step.
This is particularly important in an epidemic context where many compartment states (such as the number of currently infected but undetected individuals) cannot be directly measured in practice.
